\section{Moving jets on a two-slot bundle}
\label{sec:two-slot-jets}

We continue with the terminal fourfold \(Y\), its nef non-big canonical
divisor \(K\), and the polarizations of
\Cref{eq:small-volume-polarization}. Thus \(n=4\), the integer
\(\iota>0\) makes \(\iota K\) Cartier, and
\[
 L=L_t=r_t(K+tA),\qquad L^n\longrightarrow\epsilon^n,\qquad
 r_t\longrightarrow\infty .
\]
The positive number \(\epsilon\) is fixed throughout this section.
Retain the fixed locus \(Y^{\mathrm{vg}}\) from
\Cref{sec:moving-centers}, on which \Cref{prop:subvarieties} holds
and every integral curve has \(K\)-degree at least \(1/\iota\).
In particular,
\begin{equation}\label{eq:two-slot-curve-lower-bound}
 L\cdot C\ge r_t/\iota
 \quad\text{if \(C\) passes through a point of \(Y^{\mathrm{vg}}\).}
\end{equation}
A very general locus here means the complement of a countable union
of proper closed subsets. We may enlarge that union whenever finitely
or countably many additional conditions are imposed.

The scalar estimate of \Cref{prop:scalar-orders} limits vanishing
on one copy of \(Y\). We now seek the complementary estimate: a
polarization on a projective bundle over \(Y\times Y\) generates
many jets at a very general point. Its Seshadri constant measures
this local positivity \cite[Chapter~5]{Lazarsfeld2004I}.

Failure of the estimate would produce a moving base component.
Positive-dimensional fibers over either copy of \(Y\) are excluded
on a bundle slice, using general type and, for a multisection,
\Cref{prop:signed-alternative}. If both projections are generically
finite, the required contradiction has a different source:
ramification bounds reduce the family to finitely many fixed
covers, and their birational automorphisms cannot sweep the product.
We isolate that argument before applying the moving-center method.

\subsection{The bundle and its volumes}

Put \(S_0=\iota K\), an actual Cartier divisor on \(Y\). On
\(Y\times Y\), a subscript denotes pullback from the corresponding
factor. Let
\begin{equation}\label{eq:two-slot-bundle}
 \pi:Z=\PP\bigl(\cO(S_{0,1})\oplus\cO(S_{0,2})\bigr)
       \longrightarrow Y\times Y,\qquad
 \xi=c_1(\cO_Z(1)).
\end{equation}
We use the convention
\[
 \pi_*\cO_Z(a\xi)
 =\operatorname{Sym}^a\bigl(\cO(S_{0,1})\oplus\cO(S_{0,2})\bigr)
 \quad(a\ge0).
\]
We suppress pullback symbols for divisors on the base. For a positive
integer \(q\), define
\begin{equation}\label{eq:two-slot-polarization}
 P=L_1+L_2+q\xi,\qquad d=\dim Z=2n+1=9.
\end{equation}
This \(q\) is a bundle weight; it is unrelated to the irregularity
\(q(Y)=0\).

The two summands determine two disjoint sections, called the axes.
The complement of the axes is the torus open of \(Z\). Fixing either
base coordinate at a point and trivializing the constant line gives
a slice
\begin{equation}\label{eq:two-slot-slice}
 \pi_{\mathrm{sl}}:Z_{\mathrm{sl}}
 =\PP\bigl(\cO_Y\oplus\cO_Y(S_0)\bigr)\longrightarrow Y,
 \qquad P_{\mathrm{sl}}=L+q\xi_{\mathrm{sl}}.
\end{equation}
The restriction of \(P\) is identified with \(P_{\mathrm{sl}}\);
the remaining constant factor is trivialized when restricting
sections.

\begin{lemma}\label[lemma]{lem:two-slot-ambient}
The varieties \(Y\times Y\), \(Z\), and \(Z_{\mathrm{sl}}\) are
projective, \(\Q\)-factorial, and klt. The tautological classes of
the two bundles are nef, and \(P\) and \(P_{\mathrm{sl}}\) are
ample rational divisors.
Writing \(K_\Sigma=K_1+K_2\) on the product and \(K_\Sigma=K\) on a
slice, one has
\begin{equation}\label{eq:two-slot-canonical}
 K_{\mathrm{bundle}}
 =-2\xi+(\iota+1)K_\Sigma
 \le\frac{\iota+1}{r_t}\,P_{\mathrm{bundle}}
\end{equation}
in nef order.
\end{lemma}

\begin{proof}
Let \(p:\widetilde Y\to Y\) be a smooth projective resolution.
Rational singularities and the irregularity conclusion in
\Cref{prop:reduction} give
\(q(\widetilde Y)=0\). The product description of the Picard group,
or the seesaw principle and the vanishing of
\(\Hom(\Alb(\widetilde Y),\Pic^0(\widetilde Y))\), gives
\[
 \Pic(\widetilde Y\times\widetilde Y)
 =\operatorname{pr}_1^*\Pic(\widetilde Y)
  \oplus\operatorname{pr}_2^*\Pic(\widetilde Y).
\]
For a prime Weil divisor on \(Y\times Y\), take its strict transform
on this smooth product. Its Cartier class is a sum of pullbacks from
the two factors. Pushing down the corresponding linear equivalence
expresses the original Weil divisor class as a sum of pullbacks of
Weil divisor classes on \(Y\). These are \(\Q\)-Cartier because
\(Y\) is \(\Q\)-factorial. Thus \(Y\times Y\) is \(\Q\)-factorial.
Exceptional components have images of codimension at least two and
do not affect this divisor-class calculation.

The bundle over \(\widetilde Y\times\widetilde Y\) is a smooth
resolution of \(Z\). Its Picard group is the Picard group of the base
plus the integral tautological class. Pushing down as above proves
\(\Q\)-factoriality of \(Z\). The identical argument over
\(\widetilde Y\) proves it for the slice. A product of canonical
varieties is canonical, as follows directly from the product of
resolutions and the canonical discrepancy formula. Projective
bundles are smooth over their bases, so the bundles are klt as well.

The direct sum of nef line bundles is nef; hence its tautological
class is nef. That class is also relatively ample. Adding a positive
pullback of an ample class therefore makes it ample: for example,
write the sum as a positive multiple of a relatively ample class
made ample by a sufficiently large base twist, plus nef classes.
This proves the assertions about \(P\) and \(P_{\mathrm{sl}}\).
Finally, the canonical bundle formula for a rank-two projective
bundle gives the equality in \Cref{eq:two-slot-canonical}.
Since \(\xi\) is nef and \(K_\Sigma\le L_\Sigma/r_t\), the stated
nef inequality follows.
\end{proof}

\begin{lemma}\label[lemma]{lem:two-slot-volumes}
Assume \(q\iota/r_t\le1\). For sufficiently small \(t\), there are
positive constants \(c_n,C_n\), depending only on \(n\), such that
\begin{equation}\label{eq:two-slot-total-volume}
 c_n q\epsilon^{2n}\le P^d\le C_n q\epsilon^{2n}.
\end{equation}
On either slice,
\begin{equation}\label{eq:two-slot-slice-volume}
 0<P_{\mathrm{sl}}^{\,n+1}
 \le2^{n+2}q\epsilon^n .
\end{equation}
\end{lemma}

\begin{proof}
For a rank-two split bundle with first Chern classes \(E_1,E_2\),
the projective-bundle formula is
\[
 \pi_*(\xi^i)=\sum_{a+b=i-1}E_1^aE_2^b\quad(i\ge1),
 \qquad \pi_*(1)=0 .
\]
All classes in the expansion of \(P^d\) are nef. Put
\(L_\Sigma=L_1+L_2\). Each \(E_j=S_{0,j}\) is bounded above in nef
order by \((\iota/r_t)L_\Sigma\). Consequently
\begin{align*}
 d qL_\Sigma^{2n}
 &\le P^d\\
 &\le
 qL_\Sigma^{2n}
 \sum_{i=1}^d i\binom di(q\iota/r_t)^{i-1}
 \le d2^{d-1}qL_\Sigma^{2n}.
\end{align*}
Here \(L_\Sigma^{2n}=\binom{2n}{n}(L^n)^2\), and
\(L^n\to\epsilon^n>0\). This proves
\Cref{eq:two-slot-total-volume}.

On a slice one summand is trivial, so
\[
 P_{\mathrm{sl}}^{\,n+1}
 =\sum_{i=1}^{n+1}\binom{n+1}{i}
 q^i L^{n+1-i}S_0^{i-1}
 \le qL^n\sum_{i=1}^{n+1}\binom{n+1}{i}.
\]
For small \(t\), we may use \(L^n\le2\epsilon^n\), giving
\Cref{eq:two-slot-slice-volume}. Positivity follows either from
ampleness or from the \(i=1\) term.
\end{proof}

We now choose the parameters in the order required by the proof.
After fixing \(\epsilon\), choose \(\delta>0\) such that
\begin{equation}\label{eq:two-slot-gap-choice}
 (2/\delta)^n2^{n+2}\epsilon^n<1 .
\end{equation}
Next choose an integer \(q>0\) so large that the lower bound in
\Cref{eq:two-slot-total-volume} gives
\begin{equation}\label{eq:two-slot-jet-volume}
 P^d>(2n+3+d\delta)^d
\end{equation}
for all sufficiently small \(t\). This is possible because that
lower bound is linear in \(q\). The integer \(q\) is now fixed.
Finally decrease \(t\), imposing \(q\iota/r_t\le1\) and all the
preceding estimates. Further decreases of \(t\) do not change
\(\epsilon,\delta,\) or \(q\).

\subsection{The large Seshadri estimate}

For an ample rational divisor \(P\) and a smooth point \(z\), write
\[
 \varepsilon(P;z)
 =\inf_{C\ni z}\frac{P\cdot C}{\mult_z C},
\]
where the infimum is over integral projective curves through \(z\).

\begin{proposition}[Two-slot jets]\label[proposition]{prop:two-slot-jets}
Let \(Z\) and \(P\) be defined by
\Crefrange{eq:two-slot-bundle}{eq:two-slot-polarization}.
Fix \(\epsilon>0\), choose \(\delta\) satisfying
\Cref{eq:two-slot-gap-choice}, and then fix \(q\) large enough
for \Cref{eq:two-slot-jet-volume}. For all sufficiently small
positive rational \(t\), one has
\begin{equation}\label{eq:two-slot-seshadri}
 \varepsilon(P;z)>2n+2=10
\end{equation}
at a very general smooth point of the torus open of \(Z\).
\end{proposition}

\subsection{Dominant finite correspondences}

The final case of the jet argument will be a family of
\(n\)-dimensional subvarieties whose two projections to \(Y\) are
generically finite. The following proposition excludes such a family
using only bounds for its degree and canonical intersection. The
family itself may change with \(t\).

\begin{proposition}[Bounded correspondences cannot sweep the bundle]
\label[proposition]{prop:two-slot-correspondences}
Fix \(\epsilon>0\), an integer \(q>0\), and positive constants
\(B_0,B_1\). Keep \(Y,K,A,L_t\) and \(P=L_1+L_2+q\xi\) as above.
For all sufficiently small positive rational \(t\), there is no
irreducible finite-type parameter space \(T\) and integral closed
subvariety \(\Gamma\subset T\times Z\), dominant over \(T\), with
the following properties:
\begin{enumerate}[label=\textup{(\roman*)}]
\item The evaluation \(\Gamma\to Z\) is dominant, and a general
fiber \(V\) over \(T\) is an integral subvariety of dimension \(n\).
\item On a smooth projective resolution \(V^*\to V\), both
projections \(g_i:V^*\to Y\), \(i=1,2\), are dominant and
generically finite.
\item Writing \(P\) also for its pullback to \(V^*\), one has
\[
 0<P^n\cdot V\le B_0,\qquad
 K_{V^*}P^{n-1}\le B_1(P^n\cdot V).
\]
\end{enumerate}
The restriction on \(t\) depends only on the fixed geometry and
\(\epsilon,q,B_0,B_1\), not on the parameter space or family.
\end{proposition}

\begin{proof}
We first rule out branch divisors that sweep \(Y\), using the
intersection bounds uniformly in \(t\). We then fix \(t\) and the
family: a common branch complement gives finitely many covers, whose
birational graph families lead to an abelian variety dominating \(Y\).

Since \(P-g_i^*L\) is nef, the projection formula gives
\begin{equation}\label{eq:two-slot-cover-degrees}
 \deg(g_i)L^n\le P^n\cdot V\le B_0.
\end{equation}
The degrees are uniformly bounded because \(L^n\to\epsilon^n>0\).
This degree bound alone does not give a finite list of covers; we
must first control their branch loci.

\paragraph{Numerical bounds for branch divisors.}
Resolve the general members of the incidence family in family.
After shrinking the irreducible parameter space, there is a smooth
projective family \(\mathcal V^*\to T\) with integral fibers
birational to the corresponding \(V\), together with the two
generically finite morphisms \(g_i:V^*\to Y\) on each fiber.
For either projection, let \(T_i\to Y\) denote the finite normal
Stein factor. Its function field is \(\C(V)\).

If \(J\) is a branch prime of \(T_i\to Y\), a ramification prime on
\(T_i\) has a strict transform \(R\) on \(V^*\). A proper birational
morphism to a normal variety is an isomorphism at the generic point
of each target divisor, so this transform is present. Over the
smooth locus of \(Y\), its ramification coefficient is at least one.
There is an effective canonical comparison
\begin{equation}\label{eq:two-slot-ramification}
 K_{V^*}\sim_{\Q}g_i^*K+E_i,\qquad E_i\ge R .
\end{equation}
For completeness, resolve the rational lift to a resolution of the
terminal target. Ramification between smooth varieties and the
effective terminal discrepancies give this formula upstairs;
pushing down to \(V^*\) gives \Cref{eq:two-slot-ramification}.

Set \(d_J=L^{n-1}\cdot J\). Projection and nef order imply
\begin{equation}\label{eq:branch-degree}
 0<d_J
 \le P^{n-1}\cdot R
 \le K_{V^*}P^{n-1}
 \le C_3 ,
\end{equation}
where \(C_3=B_1B_0\) is independent of \(t\) and of the family.
The same estimate applies to every prime
divisorial image of ramification that meets \(Y_{\mathrm{sm}}\).

We also need a canonical bound on \(J\), which need not be normal.
On \(Y_{\mathrm{sm}}\), the divisor \(J\) is a Gorenstein
hypersurface. If \(\nu:J^\nu\to J\) is its normalization, finite
duality gives, in codimension one,
\[
 K_{J^\nu}+D_{\mathrm{cond}}
 \sim_{\Q}\nu^*((K+J)|_J),\qquad D_{\mathrm{cond}}\ge0 .
\]
The conductor therefore has the sign that decreases the
normalized canonical intersection. Terminal singularities are
smooth in codimension two, so \(J\cap Y_{\mathrm{sing}}\) has
dimension at most \(n-3\). Its inverse image on the normalization
has codimension at least two. It introduces no omitted
codimension-one term in this comparison. On a resolution \(J^*\),
the exceptional canonical cycles push to zero against \(n-2\)
pullbacks of \(L\). Consequently
\begin{align}
 K_{J^*}L^{n-2}
 &\le (K+J)J L^{n-2} \notag\\
 &\le d_J/r_t+d_J^2/L^n
 \le C_4d_J .
 \label{eq:branch-canonical}
\end{align}
Here \(K\le L/r_t\). The square term is the Hodge-index inequality
\[
 J^2L^{n-2}\le\frac{(JL^{n-1})^2}{L^n}.
\]
It applies to the \(\Q\)-Cartier divisor \(J\); one can pull back
to a resolution and approximate the nef class \(L\) by ample
classes. The bound in \Cref{eq:branch-degree} and the positive lower
bound for \(L^n\) make \(C_4\) uniform.

If a family of these prime images dominates \(Y\), choose a very
general member and then a very general point on its resolution
mapping to \(Y^{\mathrm{vg}}\). Such a choice is legitimate despite
the countable exceptional set. For each excluded proper closed
subset of \(Y\), dominance says that the generic image divisor is
not contained in it; omit the corresponding proper parameter
locus. On the remaining very general member, omit the inverse
images of the countably many excluded subsets, together with the
finite birational-system and quasi-finiteness exceptions.
The chosen image divisor is of general type by
\Cref{prop:subvarieties}. Its top \(L\)-degree is \(d_J\), and
\Cref{eq:branch-canonical} is the required canonical bound.
\Cref{lem:degree-obstruction} excludes this sweeping family for
sufficiently small \(t\).

\paragraph{A fixed branch complement and finitely many covers.}
The preceding exclusion holds uniformly over all families in the
statement. Now fix one such sufficiently small \(t\) and its incidence
family. No further decrease of \(t\) will be made in this proof.
The smooth projective family
\(\mathcal V^*\to T\) has integral fibers, and both projections
to the normal variety \(Y\) are dominant and generically finite.
Their degrees are bounded by \Cref{eq:two-slot-cover-degrees}.
After a finite parameter extension and shrinking, follow the
finitely many geometric generic components of the relative
Jacobian over \(Y_{\mathrm{sm}}\). Every followed component whose
fiber image is a divisor has proper total image closure: the
preceding argument excludes a dominant such family using
\Cref{eq:branch-degree,eq:branch-canonical}.

These are exactly the hypotheses of the simultaneous finite-cover
lemma \cite[Lemma~7.5]{OpenAIAbundance2026}. It gives one smooth
nonempty open \(Y^\circ\subset Y\) over which every general finite
Stein cover in either slot is finite \'etale, and finite lists of
the resulting normal covers of \(Y\). The lemma accounts for every
branch prime by its ramification strict transform, then applies
purity and finiteness of bounded-degree \'etale covers of the fixed
open. In particular, this is a finite-type family argument; a
numerical branch bound alone would not suffice. The open and lists
may depend on the fixed \(t\) and incidence family. No uniformity
in those choices is needed.

\paragraph{Countably many birational graph families.}
Each \(V\) gives a birational map between one cover from each list:
both finite Stein factors have function field \(\C(V)\). Its
image in \(Y\times Y\) is the finite image of the closure of that
birational graph. Fix one pair of covers that are birational.
A smooth projective resolution of their common birational class
has pseudo-effective canonical class, by the effective canonical
comparison with the nef \(K\) on \(Y\). It is non-uniruled by
\cite{BDPP2013}.

The birational-group structure theorem of Hanamura
\cite{Hanamura1988}, in the form stated in
\cite[Proposition 3.7 and Theorem 3.8]{Blanc2017}, permits a smooth
projective model \(U\) of this class for which the reduced
flat-graph birational scheme is a group scheme locally of finite
type and
\[
 A_U=\Aut^0(U)
\]
is its identity component, an abelian variety. The flat graphs are
represented in graph Hilbert schemes. There are countably many
Hilbert polynomials and finitely many components in each
finite-type Hilbert scheme, so the birational maps occur in
countably many translates \(A_U b\).

For such a translate, its graphs are parametrized on a dense open
by the rational evaluation map
\begin{equation}\label{eq:abelian-graph-evaluation}
 U\times A_U\dashrightarrow U\times U,\qquad
 (x,a)\longmapsto(x,a\cdot b(x)).
\end{equation}
If this map is not dominant, its image closure is a proper closed
subset of \(U\times U\). Fix birational identifications of \(U\)
with the two covers. Every graph meets the product of their fixed
birational opens densely, because it dominates both factors.
Transporting the image closure through those opens, then taking
its closure in the product of covers, preserves its dimension.
The subsequent finite map to \(Y\times Y\) also preserves
dimension. Thus all correspondences belonging to this translate
are contained in a proper closed subset of \(Y\times Y\).
The possible boundary of an individual birational map introduces
no new component: the generic graph has already been included.

There are finitely many cover pairs and countably many translates
for each pair. If every map in \Cref{eq:abelian-graph-evaluation}
were nondominant, the image of the original incidence would be
contained in a countable union of proper closed subsets of
\(Y\times Y\). But \(\Gamma\to Z\) is dominant, so its composite
to \(Y\times Y\) is dominant and its image contains a nonempty
open subset. Over the uncountable field \(\C\), such an open
cannot be covered by countably many proper closed subsets.
Therefore the map in \Cref{eq:abelian-graph-evaluation} is dominant for at
least one translate.

\paragraph{The abelian curve contradiction.}
For a general \(x\) in this last evaluation, \(A_U\cdot b(x)\)
is a dense orbit in \(U\). A stabilizer element fixes every point
of that orbit by commutativity; it fixes \(U\) by density.
The action of \(\Aut^0(U)\) is faithful, so the stabilizer is
zero-dimensional. Hence the orbit map from \(A_U\) is generically
finite dominant. Composing with the rational map through either
fixed cover gives a generically finite dominant rational map
\[
 A_U\dashrightarrow Y .
\]
In particular \(\dim A_U=\dim Y=n\).

Resolve this map as
\(\pi_A:T_A\to A_U\) and \(g:T_A\to Y\), with \(T_A\) smooth and
projective. Since the canonical bundle of \(A_U\) is trivial,
and since \(Y\) is terminal, the canonical comparisons give
\begin{equation}\label{eq:abelian-canonical-comparisons}
 K_{T_A}\sim_{\Q}E_{\pi_A},\qquad
 K_{T_A}\sim_{\Q}g^*K+R,\qquad
 E_{\pi_A},R\ge0 ,
\end{equation}
where \(E_{\pi_A}\) is exceptional over \(A_U\).
The images of the exceptional divisors of \(\pi_A\) have
codimension at least two on the smooth abelian variety.

Choose a point in the isomorphism open of \(\pi_A\) where \(g\)
is quasi-finite, outside \(\Supp(R)\), and mapping to
\(Y^{\mathrm{vg}}\). Dominance and the countable-exclusion
description ensure that such a point exists. A general sufficiently
ample complete-intersection curve in \(A_U\) through its image
avoids all the exceptional centers: they have codimension at
least two and do not contain the prescribed point. Its strict
transform \(C\) is disjoint from \(E_{\pi_A}\) and is not contained
in \(R\). The map \(g\) is nonconstant on \(C\), since it is
quasi-finite at the chosen point. Intersecting
\Cref{eq:abelian-canonical-comparisons} with \(C\) gives
\[
 0=K_{T_A}\cdot C=g^*K\cdot C+R\cdot C>0 .
\]
The final strict inequality uses
\(K\cdot g(C)\ge1/\iota\), the defining curve property of
\(Y^{\mathrm{vg}}\), multiplied by the positive degree of
\(C\to g(C)\), together with \(R\cdot C\ge0\).
This contradiction excludes the remaining correspondence family.
\end{proof}

\subsection{Excluding moving base components}

\begin{proof}[Proof of \Cref{prop:two-slot-jets}]
All constants used to exclude moving centers below are independent
of the section degree and of small \(t\), once
\(\epsilon,\delta,q\) have been fixed. We first produce such a center
from failure of \Cref{eq:two-slot-seshadri}; we then exclude each
possible dimension of its fibers over the two factors.

\paragraph{A moving component from failure of the estimate.}
Suppose that \Cref{eq:two-slot-seshadri} fails at very general
torus points. At such a point there is a curve whose ratio is
strictly less than \(2n+3\). This conclusion does not require the
infimum defining the Seshadri constant to be attained.

Take \(d+1\) levels
\[
 \tau_j=2n+3+j\delta,\qquad 0\le j\le d .
\]
By \Cref{eq:two-slot-jet-volume}, ample Hilbert asymptotics and the
dimension of the jet space at a smooth point show that, for large
divisible \(k\), every kernel of jets of order
\(\ceil{k\tau_j}\) is nonzero. Explicitly, the leading coefficients
of the two dimensions are \(P^d/d!\) and \(\tau_j^d/d!\).
If a section in the lowest kernel does not contain the curve just
chosen, its local intersection with that curve at the marked point
is at least
\(\ceil{k\tau_0}\mult_z C\). This exceeds \(kP\cdot C\).
Thus every section in that kernel contains the curve, and the
lowest base locus has a positive-dimensional component through the
mark.

These conditions produce a dominant marked-point family in the
sense of \Cref{lem:moving-multiplicity}. For generality, choose the
marked point outside the proper
rank-jumping and base-dimension loci for all divisible degrees; there
are only countably many such conditions. For the degree \(k\) now
chosen, restrict to the open where the finitely many jet kernels have
constant rank. The dimension of the base
scheme at its marked point is upper semicontinuous. The existence
of the witnessing curve at very general marks therefore forces the
positive-dimensional-base condition on a dominant locus; no family
of the witnessing curves is assumed. The finitely many generic
base components and their inclusions can then be followed after a
finite parameter extension and shrinking.

We may increase the divisible integer \(k\) to arrange, in addition,
that all divisors
\begin{equation}\label{eq:two-slot-fillers}
 kL+(2kq-j)S_0,\qquad 0\le j\le kq ,
\end{equation}
are Cartier and globally generated. Here is one simultaneous
verification. For \(1\le a\le n\), subtract \(aA\) and then \(K\).
The resulting class is
\[
 \bigl(kr_t+(2kq-j)\iota-1\bigr)K+(kr_tt-a)A,
\]
which is ample once \(kr_tt>n\) and \(kr_t\ge1\).
Klt Kawamata--Viehweg vanishing followed by regularity with respect
to \(A\) proves global generation for every \(j\) in
\Cref{eq:two-slot-fillers}.

The moving-multiplicity lemma supplies an irreducible parameter
space \(T\), a dominant incidence
\(\Gamma\subset T\times Z\), and integral fiber components
\(V=V_u\) of a fixed dimension \(1\le f<d\). The high-kernel
subseries has \(V\) as an isolated base component generically, with
ordinary generic order at least
\[
 h=\ceil{k\delta/2},\qquad k/h\le2/\delta .
\]
For a smooth resolution \(V^*\), the base-component estimates and
\Cref{eq:two-slot-canonical} give
\begin{align}
 0<P^f\cdot V
 &\le(2/\delta)^{d-f}P^d,
 \label{eq:two-slot-center-volume}\\
 K_{V^*}P^{f-1}
 &\le
 \left(\frac{\iota+1}{r_t}+\frac{2(d-f)}{\delta}\right)P^f\cdot V .
 \label{eq:two-slot-center-canonical}
\end{align}
The right sides are bounded uniformly in the specified sense.

We first exclude positive-dimensional slot fibers. The chosen incidence
point will ensure that all image subvarieties below meet the prescribed
very general locus of \(Y\). For an intermediate fiber \(F\) in a slice, let \(W\subseteq Y\)
be its image and write \(s=\dim F\), \(w=\dim W\). The three
possibilities are summarized below; the subsequent paragraphs prove
all the indicated exclusions.
\begin{center}
\small
\begin{tabular}{@{}p{0.18\textwidth}p{0.31\textwidth}p{0.43\textwidth}@{}}
\toprule
Dimensions & Geometry & Contradiction \\
\midrule
\(s=w<n\) & Generically finite over a proper image & General type and the degree obstruction on \(F\). \\[3pt]
\(s=w+1\), \(w<n\) & Entire bundle over \(W\) & Direct slice-volume bound if \(w=0\); fillers and subadjunction on \(W\) if \(w>0\). \\[3pt]
\(s=w=n\) & Non-axis multisection over \(Y\) & Its two axis intersections give a signed representative; the logarithmic degree obstruction applies. \\
\bottomrule
\end{tabular}
\end{center}
A full slot fiber is excluded by exchanging slots. This leaves
generically finite projections. Proper images are again excluded by
general type, and dominant images are excluded by
\Cref{prop:two-slot-correspondences}.

\paragraph{Restriction to a slot fiber.}
Consider \(V\) over either factor of \(Y\times Y\). Choose a very
general incidence point first. We require that its two base
coordinates belong to \(Y^{\mathrm{vg}}\), that it lies in the torus,
and that the ordinary ideal-power and isolated-component conclusions
above hold near it. We also impose the smoothness conditions for
the incidence over \(T\) and for \(V\) over the smooth open of its
actual slot image. These are available on dense opens.

Now take the slot value of this point, and let \(F\) be the reduced
fiber component through it. Write \(s=\dim F\), the general fiber
dimension. Suppose first that
\[
 0<s<n+1.
\]
The restricted high-kernel subseries on the opposite slice is
nonzero and has \(F\) as an isolated base component generically.
Indeed, near the selected point the original base support is just
\(V\). No other base component can contain the fiber component
through that point. The scheme fiber over the ambient slot agrees
with the scheme fiber over the actual image, and smoothness over
that image makes it reduced at the selected generic component.
Thus the image of \(\cI_V^h\) in the slice is contained in
\(\cI_F^h\), as an ordinary ideal power. A restriction that vanished
identically on the whole slice would contradict this local proper
base support.

Applying \Cref{lem:base-component} in the slice gives
\begin{align}
 0<P_{\mathrm{sl}}^{\,s}\cdot F
 &\le(2/\delta)^{n+1-s}P_{\mathrm{sl}}^{\,n+1},
 \label{eq:two-slot-fiber-volume}\\
 K_{F^*}P_{\mathrm{sl}}^{\,s-1}
 &\le
 \left(\frac{\iota+1}{r_t}+\frac{2(n+1-s)}{\delta}\right)
 P_{\mathrm{sl}}^{\,s}\cdot F .
 \label{eq:two-slot-fiber-canonical}
\end{align}
Let \(W\subseteq Y\) be the image of \(F\) in the other base, and put
\(w=\dim W\). It meets the prescribed very general locus. Since the
bundle fiber has dimension one, either \(w=s\) or \(w=s-1\).
We treat every possibility.

\paragraph{A generically finite proper image.}
If \(w=s<n\), the resolution of \(W\) is of general type by
\Cref{prop:subvarieties}. The same holds for \(F^*\): after resolving
the generically finite map to a resolution of \(W\), ramification
adds an effective divisor to the pulled-back canonical divisor.
The polarization \(P_{\mathrm{sl}}\) dominates the pullback of \(L\).
\Cref{eq:two-slot-fiber-volume,eq:two-slot-fiber-canonical} therefore contradict
\Cref{lem:degree-obstruction} for sufficiently small \(t\).

\paragraph{The whole bundle over a proper image.}
Suppose \(s=w+1\), where necessarily \(w<n\). A positive-dimensional
closed subset of the generic \(\PP^1\) fiber is the entire fiber.
Consequently
\[
 F=\pi_{\mathrm{sl}}^{-1}(W)
\]
as reduced integral subvarieties. Nef intersection and the
projective-bundle formula give
\begin{equation}\label{eq:whole-bundle-base-volume}
 qL^w\cdot W\le P_{\mathrm{sl}}^{\,w+1}\cdot F .
\end{equation}
For \(w=0\), the left side is \(q\). The right side is at most
\[
 (2/\delta)^nP_{\mathrm{sl}}^{\,n+1}
 \le(2/\delta)^n2^{n+2}q\epsilon^n<q
\]
by \Cref{eq:two-slot-gap-choice}, a contradiction.

Assume \(w>0\). We will apply the degree obstruction to \(W\), rather
than to the ruled variety \(F\). Expand the restricted high sections
using the two homogeneous fiber coordinates. Their coefficients of
weight \(j\), for \(0\le j\le kq\), are actual sections of
\[
 \cO_Y(kL+jS_0).
\]
Let \(\mathfrak a\) be the ideal generated by all these coefficients
in the regular local ring
\[
 R=\cO_{Y,\eta_W},\qquad \mathfrak m=\mathfrak m_R.
\]
This ring is regular because the generic point of \(W\) lies in the
smooth locus of \(Y\).

We need both the order and the support of this coefficient ideal.
After choosing frames, the local ring of the slice at the generic
point of \(F\) is
\[
 R[u]_{\mathfrak m R[u]} .
\]
The residue of \(u\) is transcendental over the residue field of
\(R\). If a coefficient had order less than \(h\), the least
\(\mathfrak m\)-adic homogeneous part of its polynomial would remain
nonzero after adjoining that residue variable. Therefore membership
of every high section in
\(\mathfrak m^hR[u]_{\mathfrak m R[u]}\) implies
\(\mathfrak a\subseteq\mathfrak m^h\).
Moreover, if a prime \(\mathfrak p\subsetneq\mathfrak m\) contained
\(\mathfrak a\), its extension to this local polynomial ring would
be a proper prime strictly below the generic ideal of \(F\)
containing the whole high base ideal. This contradicts the isolated
support of \(F\). Hence
\begin{equation}\label{eq:two-slot-coefficient-ideal}
 \mathfrak a\subseteq\mathfrak m^h,\qquad
 \sqrt{\mathfrak a}=\mathfrak m .
\end{equation}

Multiply every weight-\(j\) coefficient by the globally generated
space in \Cref{eq:two-slot-fillers}. The products all belong to one
linear subseries of
\[
 H^0\bigl(Y,\cO_Y(2k(L+qS_0))\bigr).
\]
At \(\eta_W\), each filler system generates a unit after choosing a
frame. Thus this new subseries has precisely the joint ideal
\(\mathfrak a\), not merely an ideal with the same support. Its
threshold construction in \Cref{lem:base-component} gives an
effective rational divisor
\[
 \Theta\sim_{\Q}c\,2(L+qS_0),\qquad
 0<c\le(n-w)k/h,
\]
such that the pair is lc at the generic point of \(W\) and has
an lc place centered there. Apply
\Cref{lem:numerical-subadjunction} with the testing class \(L\).
Since \(qS_0\le L\) in nef order, we obtain
\begin{align}
 K_{W^*}L^{w-1}
 &\le\bigl(K+2c(L+qS_0)\bigr)L^{w-1}\cdot W \notag\\
 &\le\left(\frac1{r_t}+\frac{8(n-w)}{\delta}\right)L^w\cdot W .
 \label{eq:whole-bundle-canonical}
\end{align}
The top degree \(L^w\cdot W\) is bounded by
\Cref{eq:whole-bundle-base-volume} and
\Cref{eq:two-slot-fiber-volume}. The variety \(W^*\) is of general
type. \Cref{eq:whole-bundle-canonical} and the top-degree
bound now contradict \Cref{lem:degree-obstruction} on \(W^*\).

\paragraph{A non-axis multisection over all of \(Y\).}
The only remaining intermediate-fiber case is \(s=w=n\).
Here \(F\) is a multisection of positive degree \(e\) over \(Y\).
It is not an axis, since it contains the selected torus point.
Let \(\Sigma_0,\Sigma_\infty\) be the axes on the slice, with
classes
\[
 \Sigma_0\sim\xi_{\mathrm{sl}},\qquad
 \Sigma_\infty\sim\xi_{\mathrm{sl}}-S_0.
\]
The intersections with \(F\) are effective integral cycles. Their
pushforwards define effective integral Weil divisors \(D_0,D_\infty\)
on \(Y\), allowing a divisor to be zero. Projection and rational
equivalence of cycles give
\begin{equation}\label{eq:axis-linear-equivalence}
 D_0-D_\infty\sim eS_0=e\iota K .
\end{equation}
This is linear equivalence of Weil divisors: codimension-one
rational equivalence on the normal variety \(Y\) is precisely
linear equivalence. In particular,
\[
 J=\frac{D_0-D_\infty}{e\iota}\sim_{\Q}K
\]
is an actual signed rational representative.

For either axis, \(P_{\mathrm{sl}}-q\Sigma\) is nef: it is \(L\)
or \(L+qS_0\). As \(F\cap\Sigma\) is effective and
\(P_{\mathrm{sl}}\ge\pi_{\mathrm{sl}}^*L\) in nef order,
\begin{equation}\label{eq:axis-degree-bound}
 qL^{n-1}\cdot D_i
 \le qP_{\mathrm{sl}}^{\,n-1}\cdot(F\cap\Sigma_i)
 \le P_{\mathrm{sl}}^{\,n}\cdot F,\qquad i=0,\infty .
\end{equation}
Take a log resolution \(p:Y^\dagger\to Y\) of \(D_0+D_\infty\),
and let \(G\) be the reduced snc divisor consisting of their strict
supports and all exceptional divisors. By
\Cref{prop:signed-alternative},
\(K_{Y^\dagger}+G\) is big. The exceptional terms push to zero
against \(L^{n-1}\), while the reduced strict support is bounded
by the integral divisor \(D_0+D_\infty\). Hence
\begin{equation}\label{eq:multisection-log-bound}
 (K_{Y^\dagger}+G)L^{n-1}
 \le \frac{L^n}{r_t}
       +L^{n-1}\cdot(D_0+D_\infty).
\end{equation}
The second term is uniformly bounded by
\Cref{eq:axis-degree-bound} and
\Cref{eq:two-slot-fiber-volume}. Also \(L^n\) is bounded above
and bounded away from zero. Thus
\Cref{eq:multisection-log-bound} is the canonical-to-volume bound
in the log version of \Cref{lem:degree-obstruction}, with
polarization \(p^*L\) and map \(p\). That lemma gives the
contradiction in this final intermediate-fiber case.

\paragraph{Reduction to finite correspondences.}
We have excluded all slot-fiber dimensions \(1,\ldots,n\).
A fiber dimension \(n+1\) over one slot image \(W\) would force
\(V\) to be the entire bundle over \(W\times Y\).
Since \(V\ne Z\), one has \(\dim W<n\); its fiber over the other
slot then has dimension \(\dim W+1\), which has just been excluded.
It follows that both projections of \(V\) to \(Y\) are generically
finite onto their images.

If one image were proper, it would have dimension \(f<n\) and be of
general type. Its generically finite cover \(V^*\) would also be
of general type, and the total bounds
\Crefrange{eq:two-slot-center-volume}{eq:two-slot-center-canonical}
would contradict \Cref{lem:degree-obstruction}. Therefore
\[
 f=n,\qquad g_i:V^*\longrightarrow Y\quad(i=1,2)
\]
are generically finite dominant morphisms. The bounds
\Crefrange{eq:two-slot-center-volume}{eq:two-slot-center-canonical}
supply constants \(B_0,B_1\) independent of small \(t\), of the
section degree \(k\), and of the family. We are therefore in the
situation excluded by \Cref{prop:two-slot-correspondences}.

Every possibility for a moving component has now been excluded.
The slice degree obstructions and
\Cref{prop:two-slot-correspondences} impose restrictions on \(t\)
independent of \(k\) and of the moving family.
The assumed failure of \Cref{eq:two-slot-seshadri} is therefore
impossible for all sufficiently small \(t\), as required.
\end{proof}

\subsection{One finite system of jets on a smooth model}

We now convert the strict Seshadri bound into a single surjective
jet-evaluation map on a smooth model. Choose \(t\) sufficiently small
for both \Cref{prop:scalar-orders,prop:two-slot-jets}, and keep it
fixed. The next corollary then supplies a point, a Cartier multiple,
and finitely many sections whose jets span the required target.
These are the finite data used in reduction to positive
characteristic; no family of jet systems over varying \(t\) is needed.

\begin{corollary}\label[corollary]{cor:fixed-two-slot-jets}
Fix sufficiently small \(t\) as in
\Cref{prop:scalar-orders,prop:two-slot-jets}, and let
\(p:W\to Y\) be any fixed smooth projective resolution. Continue
to denote pullbacks by \(L\) and \(S_0\). On
\[
 Z_W=\PP\bigl(\cO_W(S_0)_1\oplus\cO_W(S_0)_2\bigr)
       \longrightarrow W\times W,\qquad
 M=L_1+L_2+q\xi ,
\]
there are a smooth torus point \(z_*\), over the isomorphism opens
of \(p\) in both slots, and a sufficiently divisible integer
\(k_0>0\) such that \(k_0L\) is Cartier and
\begin{equation}\label{eq:fixed-two-slot-jet-map}
 H^0(Z_W,\cO_{Z_W}(k_0M))
 \longrightarrow
 \cO_{Z_W}(k_0M)\otimes
 \bigl(\cO_{Z_W,z_*}/\mathfrak m_{z_*}^{\,10k_0+1}\bigr)
\end{equation}
is surjective. In particular, finitely many sections can be fixed
whose jets span its target.
\end{corollary}

\begin{proof}
Choose a very general smooth torus point \(z\) of \(Z\) above the
isomorphism open of \(p\) in both slots and satisfying
\Cref{eq:two-slot-seshadri}, so that \(\varepsilon(P;z)>10\). If
\(b:\widehat Z\to Z\) is its blowup and \(E_z\) is the exceptional
divisor, the Seshadri criterion shows that the following rational divisor
is ample:
\[
 b^*P-10E_z.
\]
Indeed one may choose a rational
number strictly between \(10\) and \(\varepsilon(P;z)\), and then
use the usual ample interval below the Seshadri threshold.
Choose \(a>0\) so that \(aP\) and
\(a(b^*P-10E_z)\) are Cartier. Serre vanishing, applied to the
fixed sheaf \(\cO_{\widehat Z}(-E_z)\), gives
\[
 H^1\!\left(\widehat Z,
  \cO_{\widehat Z}(kb^*P-(10k+1)E_z)\right)=0
\]
for all sufficiently large multiples \(k\) of \(a\).
The standard local calculation for the blowup of a smooth point
identifies the pushforward of
\(\cO_{\widehat Z}(-(10k+1)E_z)\) with
\(\mathfrak m_z^{10k+1}\), with vanishing higher direct images.
The restriction sequence therefore gives surjectivity onto jets
through order \(10k\) at \(z\).

The natural morphism \(Z_W\to Z\) is an isomorphism near the lift
\(z_*\) of \(z\). Pullback preserves the sections in this jet
surjection and identifies their jets. Increase the divisible
integer \(k\) if necessary so that \(kL\) is Cartier, and call it
\(k_0\). Choosing finitely many inverse images of a basis of the
finite-dimensional jet target proves the assertion.
\end{proof}
